\documentclass[12pt, a4paper]{article}

% ── Shared preamble (src/lecturenotes.sty) ────────────────────────────────────
%   Packages, colors, boxes, \sol/\soleq and math macros all live in
%   src/*.sty. \input@path points LaTeX two folders up to find them.
\makeatletter
\def\input@path{{../../}}
\makeatother
\usepackage{lecturenotes}

% ── Local macros ──────────────────────────────────────────────────────────────
\providecommand{\bZ}{\mathbf{Z}}
\providecommand{\bdelta}{\boldsymbol{\delta}}

% ── Local TikZ styles (match S1) ──────────────────────────────────────────────
\usetikzlibrary{decorations.pathreplacing}
\tikzset{
  neuron/.style={circle, draw=mainblue, fill=boxbg,
      minimum size=0.6cm, inner sep=0pt, thick, font=\scriptsize},
  input/.style={circle, draw=gray!60, fill=gray!15,
      minimum size=0.6cm, inner sep=0pt, thick, font=\scriptsize},
  layerlabel/.style={font=\small\itshape, gray},
  brace/.style={decorate, decoration={brace, amplitude=4pt}, thick, gray},
  window/.style={draw=orange!80!black, line width=1.6pt},
  outcell/.style={draw=mainblue, line width=1.6pt},
}

% \drawgrid{<shift>}{<fill>}{<rows>} — grid of unit cells with values,
%   rows given as {a,b,c},{d,e,f},... ; row 1 at the top.
\newcommand{\drawgrid}[3]{%
  \begin{scope}[shift={#1}]
    \foreach \row [count=\r] in {#3} {
        \foreach \v [count=\c] in \row {
            \draw[gray!70, fill=#2] (\c-1,-\r+1) rectangle ++(1,-1);
            \node[font=\small] at (\c-0.5,-\r+0.5) {$\v$};
          }
      }
  \end{scope}%
}

% ── Document ──────────────────────────────────────────────────────────────────
\begin{document}

\begin{center}
  {\LARGE \textbf{Section 4.2: The Convolution Operation}} \\[0.3em]
  {\large\itshape Kernels, Padding, Stride, Channels and Gradients} \\[0.5em]
  {\large AI-IM Class --- Lecture Notes}%
  \ifstudentmode
    \quad {\normalsize\color{gray}[Student Copy]}%
  \fi
\end{center}

\vspace{1em}
\hrule
\vspace{1.5em}

In Section 4.1 we slid a small kernel along a signal and saw why this beats
a fully connected layer for images. In this section we make the operation
precise and compute it by hand. We
\begin{itemize}[nosep]
  \item define convolution in one and two dimensions,
  \item control the output size with \emph{padding} and \emph{stride},
  \item handle colour images and many filters (\emph{channels}),
  \item write convolution as a matrix, connecting it back to the MLP, and
  \item compute the gradient of the loss with respect to the kernel.
\end{itemize}

\paragraph{Notation.}
As in Section 3.4, bold capitals are arrays and lower-case letters are their
entries. The input is $\bX$ with entries $x_{i,j}$; the kernel is $\bW$ with
entries $w_{u,v}$; the output (pre-activation) is $\bZ$ with entries
$z_{i,j}$; $b$ is the bias. The kernel has size $k \times k$. Do not confuse
the bold kernel $\bW$ with the italic image width $W$.

% ══════════════════════════════════════════════════════════════════════════════
\section{Convolution in One Dimension}

Let $\bx = (x_1, \dots, x_n)$ be an input of length $n$ and
$\bw = (w_1, \dots, w_k)$ a kernel of length $k \le n$. Place the kernel at
position $i$, multiply the overlapping entries, and add:

\begin{resultbox}
  \[
    z_i = \sum_{u=1}^{k} w_u\, x_{i+u-1} + b,
    \qquad i = 1, 2, \dots, n - k + 1.
  \]
  The output has length $n - k + 1$.
\end{resultbox}

\begin{figure}[H]
  \centering
  \begin{tikzpicture}[x=0.9cm, y=0.9cm, >=Stealth]
    \drawgrid{(0,0)}{white}{{1,3,2,0,4}}
    \node[anchor=east, font=\small] at (-0.2,-0.5) {$\bx$};
    \draw[window] (1,0) rectangle (4,-1);

    \drawgrid{(1,-1.8)}{notebg}{{1,0,-1}}
    \node[anchor=east, font=\small] at (0.8,-2.3) {$\bw$};
    \foreach \u in {1,2,3} \draw[->, gray] (\u+0.5,-1.75) -- (\u+0.5,-1.05);

    \drawgrid{(1,-3.6)}{boxbg}{{z_1,z_2,z_3}}
    \node[anchor=east, font=\small] at (0.8,-4.1) {$\bz$};
    \draw[outcell] (2,-3.6) rectangle (3,-4.6);

    \node[font=\small, anchor=west, align=left] at (5.6,-2.1)
    {position $i = 2$:\\
      $z_2 = (1)(3) + (0)(2) + (-1)(0)$};
  \end{tikzpicture}
  \caption{One-dimensional convolution with $n = 5$ and $k = 3$ (bias $b = 0$).
    The kernel slides one step at a time; the window at position $i = 2$ is
    outlined.}
  \label{fig:conv1d}
\end{figure}

\begin{stepbox}{Exercise 1 \quad A 1-D convolution}
  Use $\bx = (1, 3, 2, 0, 4)$, $\bw = (1, 0, -1)$ and $b = 0$, as in
  Figure~\ref{fig:conv1d}.
  \begin{enumerate}[label=(\alph*), nosep]
    \item The output has length $\sol[3cm]{5 - 3 + 1 = 3}$.
    \item $z_1 = \sol[4cm]{1 - 2 = -1}$, \quad
          $z_2 = \sol[3cm]{3 - 0 = 3}$, \quad
          $z_3 = \sol[3cm]{2 - 4 = -2}$.
    \item This kernel computes $z_i = x_i - x_{i+2}$. It measures how much the
          signal $\sol[2.5cm]{\text{decreases}}$ over two steps.
  \end{enumerate}
\end{stepbox}

\begin{notebox}
  \textbf{Convolution or cross-correlation?}
  In mathematics and signal processing, \emph{convolution} flips the kernel
  before sliding it:
  $(x * w)_i = \sum_{u} x_{i-u}\, w_u$. The operation in the box above, which
  does not flip, is strictly called \textbf{cross-correlation}.
  With the flipped kernel $(-1, 0, 1)$, Exercise 1 would give
  $\bz = (1, -3, 2)$: the same numbers with opposite sign.

  In a CNN the kernel values are \emph{learned}. If the flipped version is
  better, training simply learns the flipped numbers, so the choice makes no
  difference. Deep-learning libraries, including PyTorch's
  \texttt{nn.Conv1d} and \texttt{nn.Conv2d}, compute cross-correlation and
  call it convolution. We do the same.
\end{notebox}

% ══════════════════════════════════════════════════════════════════════════════
\section{Convolution in Two Dimensions}

For an image the kernel is a small $k \times k$ grid that slides across both
rows and columns. At each position we multiply the kernel entrywise with the
$k \times k$ patch beneath it and add everything up.

\begin{resultbox}
  \[
    z_{i,j} = \sum_{u=1}^{k} \sum_{v=1}^{k} w_{u,v}\; x_{i+u-1,\; j+v-1} + b .
  \]
  For an $H \times W$ input the output is
  $(H - k + 1) \times (W - k + 1)$.
\end{resultbox}

The output is called a \textbf{feature map}: entry $z_{i,j}$ says how
strongly the pattern in the kernel is present in the patch whose top-left
corner is pixel $(i, j)$.

\begin{figure}[H]
  \centering
  \begin{tikzpicture}[x=0.72cm, y=0.72cm, >=Stealth]
    \drawgrid{(0,0)}{white}{{0,0,0,0,0},{0,0,0,0,0},{0,0,1,1,1},{0,0,1,1,1},{0,0,1,1,1}}
    \fill[orange!20, opacity=0.6] (0,0) rectangle (3,-3);
    \draw[window] (0,0) rectangle (3,-3);
    \node[layerlabel] at (2.5,-5.6) {input $\bX$ ($5\times5$)};

    \node[font=\Large] at (6.0,-2.5) {$\ast$};

    \drawgrid{(7,-1)}{notebg}{{-1,0,1},{-1,0,1},{-1,0,1}}
    \node[layerlabel, align=center] at (8.5,-4.9) {kernel $\bW_{\!v}$\\($3\times3$)};

    \node[font=\Large] at (11.0,-2.5) {$=$};

    \drawgrid{(12,-1)}{boxbg}{{z_{1,1},z_{1,2},z_{1,3}},{z_{2,1},z_{2,2},z_{2,3}},{z_{3,1},z_{3,2},z_{3,3}}}
    \draw[outcell] (12,-1) rectangle (13,-2);
    \node[layerlabel, align=center] at (13.5,-4.9) {feature map $\bZ$\\($3\times3$)};

    \draw[->, orange!80!black, thick] (3,0) to[bend left=20] (12.5,-0.95);
  \end{tikzpicture}
  \caption{Two-dimensional convolution. The input contains a bright square in
    its lower-right corner. The outlined $3\times3$ patch, multiplied entrywise
    by the kernel and summed, gives $z_{1,1}$. Sliding the patch one column
    right gives $z_{1,2}$, and so on.}
  \label{fig:conv2d}
\end{figure}

\begin{stepbox}{Exercise 2 \quad Detecting edges by hand}
  Use the input $\bX$ of Figure~\ref{fig:conv2d} with $b = 0$ and two
  kernels:
  \[
    \bW_{\!v} =
    \begin{pmatrix} -1 & 0 & 1 \\ -1 & 0 & 1 \\ -1 & 0 & 1 \end{pmatrix},
    \qquad
    \bW_{\!h} =
    \begin{pmatrix} -1 & -1 & -1 \\ 0 & 0 & 0 \\ 1 & 1 & 1 \end{pmatrix}.
  \]
  $\bW_{\!v}$ computes (right column of the patch) $-$ (left column);
  $\bW_{\!h}$ computes (bottom row) $-$ (top row).
  \[
    \bZ_v =
    \begin{pmatrix}
      \sol[0.6cm]{1} & \sol[0.6cm]{1} & \sol[0.6cm]{0} \\[2pt]
      \sol[0.6cm]{2} & \sol[0.6cm]{2} & \sol[0.6cm]{0} \\[2pt]
      \sol[0.6cm]{3} & \sol[0.6cm]{3} & \sol[0.6cm]{0}
    \end{pmatrix},
    \qquad
    \bZ_h =
    \begin{pmatrix}
      \sol[0.6cm]{1} & \sol[0.6cm]{2} & \sol[0.6cm]{3} \\[2pt]
      \sol[0.6cm]{1} & \sol[0.6cm]{2} & \sol[0.6cm]{3} \\[2pt]
      \sol[0.6cm]{0} & \sol[0.6cm]{0} & \sol[0.6cm]{0}
    \end{pmatrix}
  \]
  Where in the image is $\bZ_v$ largest, and why?
  \soleq[1.2cm]{\parbox{0.85\linewidth}{\raggedright Bottom-left: those
      patches contain the whole vertical edge of the square (dark on the left,
      bright on the right, in all three rows).}}
\end{stepbox}

The two kernels respond to different features of the \emph{same} image:
$\bW_{\!v}$ to the left (vertical) edge of the square, $\bW_{\!h}$ to its
top (horizontal) edge. Inside the square, where the image is flat, both are
zero.

\begin{notebox}
  \textbf{Hand-designed vs.\ learned kernels.}
  $\bW_{\!v}$ and $\bW_{\!h}$ are classical image-processing filters (the
  \emph{Prewitt} filters), designed by hand. In a CNN the nine numbers in each
  kernel are \emph{parameters}, found by gradient descent like any other
  weight. When we look at the kernels a trained CNN learned in its first
  layer, many turn out to be edge and colour detectors very much like these.
\end{notebox}

% ══════════════════════════════════════════════════════════════════════════════
\section{Padding}

Without help, every convolution shrinks the input: a $3 \times 3$ kernel
removes one pixel from each border, so a $28 \times 28$ image becomes
$26 \times 26$, then $24 \times 24$, and so on. Pixels on the border are also
covered by fewer windows than pixels in the middle, so the network sees them
less.

The fix is \textbf{padding}: surround the input with $p$ rows and columns of
zeros before convolving.

\begin{figure}[H]
  \centering
  \begin{tikzpicture}[x=0.6cm, y=0.6cm]
    \foreach \r in {1,...,6} \foreach \c in {1,...,6} {
          \pgfmathtruncatemacro{\border}{(\r==1)||(\r==6)||(\c==1)||(\c==6)}
          \ifnum\border=1
            \draw[gray!60, fill=gray!15] (\c-1,-\r+1) rectangle ++(1,-1);
            \node[font=\scriptsize, gray] at (\c-0.5,-\r+0.5) {0};
          \else
            \draw[gray!70, fill=white] (\c-1,-\r+1) rectangle ++(1,-1);
            \node[font=\scriptsize] at (\c-0.5,-\r+0.5) {$x$};
          \fi
        }
    \draw[thick] (1,-1) rectangle (5,-5);
    \draw[window] (0,0) rectangle (3,-3);
    \node[layerlabel, align=center] at (3,-6.9)
    {$4\times4$ input, $p = 1$\\(padded to $6\times6$)};

    \draw[->, thick, >=Stealth] (6.8,-3) -- node[above, font=\small]
    {$3\times3$ kernel} (10.2,-3);

    \foreach \r in {1,...,4} \foreach \c in {1,...,4}
    \draw[gray!70, fill=boxbg] (\c+10.5,-\r) rectangle ++(1,-1);
    \draw[outcell] (11.5,-1) rectangle (12.5,-2);
    \node[layerlabel, align=center] at (13.5,-6.9) {$4\times4$ output\\(same size)};
  \end{tikzpicture}
  \caption{Zero padding with $p = 1$. The first window (orange) is centred on
    the corner pixel, so the corner now gets its own output, and the output
    has the same size as the input.}
  \label{fig:padding}
\end{figure}

Padding adds $2p$ to each dimension, so the output length becomes
$n + 2p - k + 1$. Three choices of $p$ have names.

\begin{center}
  \renewcommand{\arraystretch}{1.4}
  \begin{tabular}{llll}
    \toprule
    Type  & Padding $p$          & Output length & Use                       \\
    \midrule
    Valid & $0$                  & $n - k + 1$   & shrinks; common           \\
    Same  & $(k-1)/2$ ($k$ odd)  & $n$           & keeps size; most common   \\
    Full  & $k - 1$              & $n + k - 1$   & grows; rare in CNNs       \\
    \bottomrule
  \end{tabular}
\end{center}

Same padding is the reason almost all CNN kernels have \emph{odd} size
($3\times3$, $5\times5$, $7\times7$): then $(k-1)/2$ is a whole number and the
kernel has a centre pixel. In PyTorch, \texttt{padding=0} is valid padding,
and \texttt{padding=1} with a $3 \times 3$ kernel (or
\texttt{padding="same"}) is same padding.

\begin{stepbox}{Exercise 3 \quad Same padding in 1-D}
  Repeat Exercise 1 ($\bx = (1, 3, 2, 0, 4)$, $\bw = (1, 0, -1)$) with same
  padding.
  \begin{enumerate}[label=(\alph*), nosep]
    \item $p = \sol[1cm]{1}$; the padded input is
          $\sol[5cm]{(0,\ 1,\ 3,\ 2,\ 0,\ 4,\ 0)}$.
    \item $\bz = \bigl(\,\sol[0.9cm]{-3},\ \sol[0.9cm]{-1},\ \sol[0.9cm]{3},\
            \sol[0.9cm]{-2},\ \sol[0.9cm]{0}\,\bigr)$, of length
          $\sol[1cm]{5}$.
    \item Compare with Exercise 1: the middle three outputs are
          $\sol[5cm]{\text{the same values } (-1, 3, -2)}$; padding only adds
          outputs at the borders.
  \end{enumerate}
\end{stepbox}

% ══════════════════════════════════════════════════════════════════════════════
\section{Stride}

So far the kernel has moved one step at a time. With \textbf{stride} $s$ it
moves $s$ steps, so only every $s$-th output is computed. A stride of $2$
roughly halves the output in each dimension. This reduces computation and,
as noted in Section 4.1, makes the receptive field of later layers grow
faster.

\begin{figure}[H]
  \centering
  \begin{tikzpicture}[x=0.85cm, y=0.85cm]
    \drawgrid{(0,0)}{white}{{0,1,3,2,0,4,0}}
    \foreach \a in {0,6} {
        \filldraw[fill=gray!15, draw=gray!70] (\a,0) rectangle ++(1,-1);
        \node[font=\small, gray] at (\a+0.5,-0.5) {$0$};
      }
    \foreach \a/\col/\h in {0/orange!80!black/0.25, 2/mainblue/0.55,
        4/green!50!black/0.85} {
        \draw[\col, line width=1.4pt] (\a+0.08, \h) -- (\a+0.08, \h+0.15)
        -- (\a+2.92, \h+0.15) -- (\a+2.92, \h);
      }
    \node[layerlabel, anchor=west] at (7.3,-0.5) {padded input ($p = 1$)};
    \node[font=\small, anchor=west] at (7.3,0.6) {windows start at $1, 3, 5$};

    \drawgrid{(2,-1.8)}{boxbg}{{z_1,z_2,z_3}}
    \node[layerlabel, anchor=west] at (7.3,-2.3) {output with $s = 2$};
  \end{tikzpicture}
  \caption{Stride $s = 2$: the kernel of width $3$ jumps two positions at a
    time, giving three outputs instead of five.}
  \label{fig:stride}
\end{figure}

Combining padding and stride gives the general output-size formula. In two
dimensions it is applied to the height and the width separately.

\begin{resultbox}
  For input size $n$, kernel size $k$, padding $p$ and stride $s$:
  \[
    n_{\text{out}} = \left\lfloor \frac{n + 2p - k}{s} \right\rfloor + 1 .
  \]
  The floor $\lfloor\cdot\rfloor$ rounds down: a window that would stick out
  past the padded input is dropped.
\end{resultbox}

\begin{stepbox}{Exercise 4 \quad Output sizes}
  \begin{enumerate}[label=(\alph*), nosep]
    \item Exercise 3 with stride $s = 2$:
          $n_{\text{out}} = \sol[4cm]{\lfloor (5 + 2 - 3)/2 \rfloor + 1 = 3}$,
          and $\bz = \bigl(\,\sol[0.9cm]{-3},\ \sol[0.9cm]{3},\
            \sol[0.9cm]{0}\,\bigr)$.
    \item Complete the table.
  \end{enumerate}
  \begin{center}
    \renewcommand{\arraystretch}{1.4}
    \begin{tabular}{ccccc}
      \toprule
      Input $n$ & Kernel $k$ & Padding $p$ & Stride $s$ & Output $n_{\text{out}}$ \\
      \midrule
      10        & 5          & 2           & 1          & $\sol[1.2cm]{10}$       \\
      10        & 3          & 2           & 2          & $\sol[1.2cm]{6}$        \\
      28        & 5          & 0           & 1          & $\sol[1.2cm]{24}$       \\
      32        & 3          & 1           & 1          & $\sol[1.2cm]{32}$       \\
      7         & 3          & 0           & 2          & $\sol[1.2cm]{3}$        \\
      224       & 7          & 3           & 2          & $\sol[1.2cm]{112}$      \\
      \bottomrule
    \end{tabular}
  \end{center}
  The last row is the first layer of ResNet (Section 4.4).
\end{stepbox}

% ══════════════════════════════════════════════════════════════════════════════
\section{Channels}

% ──────────────────────────────────────────────────────────────────────────────
\subsection{Multiple Input Channels}

A colour image has $C_{\text{in}} = 3$ channels. A kernel for it is not a
flat $k \times k$ grid but a \emph{stack} of $C_{\text{in}}$ grids, one per
channel, so its shape is $(C_{\text{in}}, k, k)$. Each slice is convolved with
its own channel, and the results are \textbf{added} together with one bias.
The kernel slides only across height and width, never across channels, so the
output is a \emph{single} feature map.

\begin{resultbox}
  \[
    z_{i,j} = \sum_{c=1}^{C_{\text{in}}} \sum_{u=1}^{k} \sum_{v=1}^{k}
    w_{c,u,v}\; x_{c,\; i+u-1,\; j+v-1} + b .
  \]
\end{resultbox}

\begin{figure}[H]
  \centering
  \begin{tikzpicture}[x=0.5cm, y=0.5cm, >=Stealth]
    % input channels
    \foreach \col/\s in {blue/2, green!60!black/1, red/0} {
        \filldraw[fill=\col!15, draw=\col!80!black, thick]
        (0.6*\s, 0.6*\s) rectangle ++(5,-5);
      }
    \filldraw[fill=orange!30, draw=orange!80!black, thick] (0,0) rectangle (2,-2);
    \node[layerlabel, align=center] at (3,-6.4) {input\\$(C_{\text{in}}, H, W)$};

    % kernel slices
    \begin{scope}[xshift=4.6cm, yshift=-0.5cm]
      \foreach \col/\s in {blue/2, green!60!black/1, red/0} {
          \filldraw[fill=\col!35, draw=\col!80!black, thick]
          (0.6*\s, 0.6*\s) rectangle ++(2,-2);
        }
      \node[layerlabel, align=center] at (1.6,-5.4)
      {one filter\\$(C_{\text{in}}, k, k)$};
    \end{scope}
    \node[font=\Large] at (8.6,-1.8) {$\ast$};

    % sum
    \node[circle, draw=mainblue, thick, minimum size=0.8cm] (S) at (18,-1.8)
    {$\Sigma$};
    \node[font=\small] at (18,-0.4) {$+\,b$};
    \draw[->, thick] (13.2,-1.8) -- node[above, font=\scriptsize, align=center]
    {multiply \& add\\per channel} (S);

    % output
    \filldraw[fill=boxbg, draw=mainblue, thick] (20.5,0.2) rectangle ++(4,-4);
    \fill[mainblue!50] (20.5,0.2) rectangle ++(1,-1);
    \draw[->, thick] (S) -- (20.5,-1.8);
    \node[layerlabel, align=center] at (22.5,-4.9) {one feature map\\$(H', W')$};
  \end{tikzpicture}
  \caption{A filter has one $k \times k$ slice per input channel. Each slice
    meets its own channel; all products are added, plus one bias, to give one
    entry of a single output feature map.}
  \label{fig:channels}
\end{figure}

\begin{stepbox}{Exercise 5 \quad Two input channels}
  An input has two channels and a filter has two $2 \times 2$ slices, with
  bias $b = 1$ and no padding:
  \[
    \bX_1 = \begin{pmatrix} 1 & 0 & 2 \\ 3 & 1 & 0 \\ 0 & 2 & 1 \end{pmatrix},\;
    \bX_2 = \begin{pmatrix} 0 & 1 & 1 \\ 2 & 0 & 1 \\ 1 & 1 & 0 \end{pmatrix},
    \qquad
    \bW_1 = \begin{pmatrix} 1 & -1 \\ 0 & 2 \end{pmatrix},\;
    \bW_2 = \begin{pmatrix} 0 & 1 \\ 1 & 0 \end{pmatrix}.
  \]
  Compute each channel's contribution, then add them and the bias.
  \[
    \underbrace{\begin{pmatrix}
        \sol[0.6cm]{3} & \sol[0.6cm]{-2} \\[2pt]
        \sol[0.6cm]{6} & \sol[0.6cm]{3}
      \end{pmatrix}}_{\bX_1 \ast \bW_1}
    +
    \underbrace{\begin{pmatrix}
        \sol[0.6cm]{3} & \sol[0.6cm]{1} \\[2pt]
        \sol[0.6cm]{1} & \sol[0.6cm]{2}
      \end{pmatrix}}_{\bX_2 \ast \bW_2}
    + 1
    =
    \underbrace{\begin{pmatrix}
        \sol[0.6cm]{7} & \sol[0.6cm]{0} \\[2pt]
        \sol[0.6cm]{8} & \sol[0.6cm]{6}
      \end{pmatrix}}_{\bZ}
  \]
\end{stepbox}

% ──────────────────────────────────────────────────────────────────────────────
\subsection{Multiple Filters}

One filter detects one kind of feature. A convolutional layer uses
$C_{\text{out}}$ filters, each with its own weights and bias, and stacks the
$C_{\text{out}}$ feature maps into an output with $C_{\text{out}}$
\emph{channels}. That output is the input of the next layer, which is why its
filters need $C_{\text{in}}$ equal to the previous layer's $C_{\text{out}}$.

\begin{figure}[H]
  \centering
  \begin{tikzpicture}[x=0.45cm, y=0.45cm, >=Stealth]
    \foreach \s in {2,1,0}
    \filldraw[fill=gray!12, draw=gray!70!black, thick]
    (0.6*\s, 0.6*\s) rectangle ++(5,-5);
    \node[layerlabel, align=center] at (3,-6.6) {input\\$(3, H, W)$};

    \foreach \f/\col/\y in {1/red/2.6, 2/mainblue/-1.0, 3/green!50!black/-4.6} {
        \begin{scope}[xshift=4.6cm, yshift=\y*0.45cm]
          \foreach \s in {2,1,0}
          \filldraw[fill=\col!25, draw=\col!80!black]
          (0.4*\s, 0.4*\s) rectangle ++(1.6,-1.6);
        \end{scope}
        \draw[->, \col!80!black, thick] (13.4, \y-0.6) -- (19.6, -0.6-0.5*\f);
      }
    \node[layerlabel, align=center] at (11.7,-8.3)
    {$C_{\text{out}} = 3$ filters, each $(3, k, k)$};

    \foreach \col/\s in {green!50!black/2, mainblue/1, red/0}
    \filldraw[fill=\col!15, draw=\col!80!black, thick]
    (20+0.6*\s, 0.6*\s) rectangle ++(4,-4);
    \node[layerlabel, align=center] at (22.6,-6.6) {output\\$(3, H', W')$};
  \end{tikzpicture}
  \caption{A layer with $C_{\text{out}} = 3$ filters produces an output with
    three channels: one feature map per filter.}
  \label{fig:filters}
\end{figure}

\begin{resultbox}
  A convolutional layer with $C_{\text{in}}$ input channels, $C_{\text{out}}$
  filters and kernel size $k \times k$ has
  \begin{itemize}[nosep]
    \item a weight tensor of shape $(C_{\text{out}}, C_{\text{in}}, k, k)$,
    \item a bias vector of shape $(C_{\text{out}})$,
  \end{itemize}
  \[
    \#\text{parameters} = (C_{\text{in}} \cdot k \cdot k + 1)\; C_{\text{out}} .
  \]
  The count does \emph{not} depend on the height or width of the image.
\end{resultbox}

Compare this with the MLP formula $(n_{\ell-1} + 1)\,n_\ell$: one filter
plays the role of one node, and its $C_{\text{in}} k^2$ weights are its
inputs. The difference is that the same filter is reused at every position.

\begin{codebox}
import torch
import torch.nn as nn

conv = nn.Conv2d(in_channels=3, out_channels=32, kernel_size=5)
x = torch.randn(64, 3, 32, 32)      # a batch of 64 colour images

print(conv.weight.shape)   # torch.Size([32, 3, 5, 5])
print(conv.bias.shape)     # torch.Size([32])
print(conv(x).shape)       # torch.Size([64, 32, 28, 28])
\end{codebox}

\begin{stepbox}{Exercise 6 \quad Shapes and parameters}
  A batch of colour images of size $32 \times 32$ passes through two
  convolutional layers (stride $1$).
  \begin{center}
    \renewcommand{\arraystretch}{1.4}
    \begin{tabular}{lccc}
      \toprule
      Layer                          & Weight shape                    & Parameters               & Output shape               \\
      \midrule
      Input                          & ---                             & ---                      & $(3, 32, 32)$              \\
      32 filters $5\times5$, valid   & $\sol[2.4cm]{(32, 3, 5, 5)}$    & $\sol[1.8cm]{2\,432}$    & $\sol[2.4cm]{(32, 28, 28)}$ \\
      64 filters $3\times3$, same    & $\sol[2.4cm]{(64, 32, 3, 3)}$   & $\sol[1.8cm]{18\,496}$   & $\sol[2.4cm]{(64, 28, 28)}$ \\
      \bottomrule
    \end{tabular}
  \end{center}
  A dense layer mapping the same input ($3 \cdot 32 \cdot 32$ values) to the
  same first output ($32 \cdot 28 \cdot 28$ values) would need
  \[
    (3072 + 1) \times 25\,088 = \sol[4cm]{77\,095\,424}
    \ \text{parameters,}
  \]
  about $\sol[2.5cm]{32\,000}$ times as many as the first convolution.
\end{stepbox}

% ══════════════════════════════════════════════════════════════════════════════
\section{Convolution Is a Constrained Dense Layer}

A convolution is linear in its input, so, like any linear map, it can be
written as a matrix--vector product $\bz = \bW \bx$, exactly as in a fully
connected layer (Section 3.4). The matrix just has a special form.

\begin{stepbox}{Exercise 7 \quad Write the convolution as a matrix}
  Take $\bx = (x_1, \dots, x_5)$ and $\bw = (w_1, w_2, w_3)$, valid
  padding and no bias. Fill in the matrix so that $\bz = \bW \bx$ gives
  $z_i = w_1 x_i + w_2 x_{i+1} + w_3 x_{i+2}$.
  \[
    \begin{pmatrix} z_1 \\ z_2 \\ z_3 \end{pmatrix}
    =
    \begin{pmatrix}
      \sol[0.6cm]{w_1} & \sol[0.6cm]{w_2} & \sol[0.6cm]{w_3} & \sol[0.6cm]{0}   & \sol[0.6cm]{0}   \\[2pt]
      \sol[0.6cm]{0}   & \sol[0.6cm]{w_1} & \sol[0.6cm]{w_2} & \sol[0.6cm]{w_3} & \sol[0.6cm]{0}   \\[2pt]
      \sol[0.6cm]{0}   & \sol[0.6cm]{0}   & \sol[0.6cm]{w_1} & \sol[0.6cm]{w_2} & \sol[0.6cm]{w_3}
    \end{pmatrix}
    \begin{pmatrix} x_1 \\ x_2 \\ x_3 \\ x_4 \\ x_5 \end{pmatrix}
  \]
\end{stepbox}

A general dense layer from $5$ inputs to $3$ outputs has $15$ free weights.
The convolution matrix has only $3$, arranged in two kinds of constraint:
\begin{itemize}[nosep]
  \item \textbf{Zeros} outside a band: each output sees only its receptive
        field (sparse interactions).
  \item \textbf{Repeated entries} along each diagonal: every row is the row
        above shifted one place right (parameter sharing).
\end{itemize}
So a convolutional layer is a fully connected layer whose weights are forced
to obey these constraints. Goodfellow et al.\ call this an
\emph{infinitely strong prior}: before seeing any data, we decide that the
useful features are local and the same everywhere. For images this belief is
correct, and that is why CNNs learn from far less data than MLPs.

% ══════════════════════════════════════════════════════════════════════════════
\section{Learning the Kernel: Backpropagation}

The kernel weights are learned by gradient descent, so we need
$\partial J / \partial w_u$. The matrix view tells us how: weight $w_u$
appears in \emph{every} row, once per output position. In Section 3.3 each
weight affected the loss along one path. Here $w_u$ affects the loss through
every $z_i$, and the chain rule adds up all the paths.

\begin{figure}[H]
  \centering
  \begin{tikzpicture}[x=1.4cm, y=1.6cm, >=Stealth]
    \foreach \i in {1,...,4} \node[input] (x\i) at (\i, 0) {$x_{\i}$};
    \foreach \j in {1,...,3} \node[neuron] (z\j) at (\j+0.5, 1) {$z_{\j}$};
    \node (J) at (2.5, 2) {$J$};
    \foreach \j in {1,...,3} {
        \pgfmathtruncatemacro{\ja}{\j+1}
        \draw[red!70!black, line width=1.3pt] (z\j) -- (x\j);
        \draw[gray!50, line width=0.6pt] (z\j) -- (x\ja);
        \draw[->, gray] (z\j) -- (J);
      }
    \node[font=\scriptsize, red!70!black, anchor=east] at (0.6, 0.5)
    {$w_1$ (red) is used three times};
    \node[font=\scriptsize, gray, anchor=west] at (4.4, 0.5)
    {$w_2$ (gray)};
    \foreach \j in {1,...,3}
    \node[font=\scriptsize, mainblue] at (\j+0.15, 1.45) {$\delta_{\j}$};
  \end{tikzpicture}
  \caption{A kernel of width $2$ on an input of length $4$. Weight $w_1$
    connects $x_1 \to z_1$, $x_2 \to z_2$ and $x_3 \to z_3$, so its gradient
    collects one term from each output.}
  \label{fig:backprop}
\end{figure}

Write $\delta_i = \partial J / \partial z_i$ for the gradient arriving at
output $i$ from later layers (as in Section 3.3). Since
$z_i = \sum_u w_u x_{i+u-1} + b$, we have
$\partial z_i / \partial w_u = x_{i+u-1}$ and $\partial z_i / \partial b = 1$.
The chain rule, summed over all outputs, gives:

\begin{resultbox}
  \[
    \pd{J}{w_u} = \sum_{i} \delta_i\; x_{i+u-1},
    \qquad
    \pd{J}{b} = \sum_{i} \delta_i .
  \]
  The gradient of a shared weight is the \textbf{sum} of the gradients from
  every position where it is used. Notice that
  $\partial J / \partial w_u$ is itself a convolution: slide $\bdelta$ along
  $\bx$.
\end{resultbox}

\begin{stepbox}{Exercise 8 \quad One gradient-descent step on a kernel}
  Let $\bx = (1, 2, -1, 3)$, $\bw = (w_1, w_2) = (1, -1)$ and $b = 0$. The
  desired output is $\by = (0, 2, -3)$, and the loss is
  $J = \tfrac{1}{2}\sum_i (z_i - y_i)^2$, so $\delta_i = z_i - y_i$.
  \begin{enumerate}[label=(\alph*)]
    \item Forward pass:
          $\bz = \bigl(\,\sol[0.9cm]{-1},\ \sol[0.9cm]{3},\ \sol[0.9cm]{-4}\,\bigr)$
          \quad and \quad
          $J = \sol[3cm]{\tfrac12(1 + 1 + 1) = 1.5}$.
    \item Output gradients:
          $\bdelta = \bigl(\,\sol[0.9cm]{-1},\ \sol[0.9cm]{1},\ \sol[0.9cm]{-1}\,\bigr)$.
    \item Kernel and bias gradients:
          \begin{align*}
            \pd{J}{w_1} & = \delta_1 x_1 + \delta_2 x_2 + \delta_3 x_3
            = \sol[5cm]{-1 + 2 + 1 = 2}                                 \\
            \pd{J}{w_2} & = \delta_1 x_2 + \delta_2 x_3 + \delta_3 x_4
            = \sol[5cm]{-2 - 1 - 3 = -6}                                \\
            \pd{J}{b}   & = \delta_1 + \delta_2 + \delta_3 = \sol[5cm]{-1}
          \end{align*}
    \item Update with learning rate $\eta = 0.1$:
          $w_1 = \sol[1.3cm]{0.8}$, \ $w_2 = \sol[1.3cm]{-0.4}$, \
          $b = \sol[1.3cm]{0.1}$.
    \item New forward pass:
          $\bz = \bigl(\,\sol[0.9cm]{0.1},\ \sol[0.9cm]{2.1},\ \sol[0.9cm]{-1.9}\,\bigr)$
          and $J = \sol[1.5cm]{0.615}$. The loss went
          $\sol[1.5cm]{\text{down}}$.
  \end{enumerate}
\end{stepbox}

PyTorch reproduces these numbers with automatic differentiation:

\begin{codebox}
conv = nn.Conv1d(in_channels=1, out_channels=1, kernel_size=2)
with torch.no_grad():
    conv.weight[:] = torch.tensor([[[1., -1.]]])
    conv.bias.zero_()

x = torch.tensor([[[1., 2., -1., 3.]]])   # shape (N, C, length)
y = torch.tensor([[[0., 2., -3.]]])
J = 0.5 * ((conv(x) - y) ** 2).sum()
J.backward()
print(conv.weight.grad)    # tensor([[[ 2., -6.]]])
print(conv.bias.grad)      # tensor([-1.])
\end{codebox}

\begin{notebox}
  \textbf{Gradient with respect to the input.}
  To keep backpropagating to earlier layers we also need
  $\partial J / \partial x_j = \sum_i \delta_i\, w_{j-i+1}$, summed over the
  outputs whose window contains $x_j$. This turns out to be a \emph{full}
  convolution of $\bdelta$ with the \emph{flipped} kernel: the flip we
  avoided earlier comes back in the backward pass. Frameworks compute it for
  us; the important idea is the one in the box above. Shared weights
  \emph{add} their gradients.
\end{notebox}

% ══════════════════════════════════════════════════════════════════════════════
\section{Summary}

\begin{center}
  \renewcommand{\arraystretch}{1.5}
  \begin{tabular}{ll}
    \toprule
    Quantity               & Formula                                                                                \\
    \midrule
    1-D convolution        & $z_i = \sum_{u} w_u\, x_{i+u-1} + b$                                                   \\
    2-D, many channels     & $z_{i,j} = \sum_{c}\sum_{u}\sum_{v} w_{c,u,v}\, x_{c,\,i+u-1,\,j+v-1} + b$              \\
    Output size            & $n_{\text{out}} = \lfloor (n + 2p - k)/s \rfloor + 1$                                 \\
    Same padding           & $p = (k-1)/2$                                                                          \\
    Parameters of a layer  & $(C_{\text{in}} k^2 + 1)\, C_{\text{out}}$                                             \\
    Kernel gradient        & $\partial J / \partial w_u = \sum_i \delta_i\, x_{i+u-1}$                              \\
    \bottomrule
  \end{tabular}
\end{center}

\begin{itemize}
  \item A convolution slides a kernel over the input and takes a weighted sum
        at each position. Deep learning uses the unflipped version
        (cross-correlation).
  \item Padding controls whether the output shrinks (valid) or keeps its size
        (same). Stride $s$ skips positions and shrinks the output about
        $s$-fold.
  \item A filter has one slice per input channel; a layer has $C_{\text{out}}$
        filters and produces $C_{\text{out}}$ feature maps. Weights in PyTorch
        have shape $(C_{\text{out}}, C_{\text{in}}, k, k)$.
  \item A convolution is a dense layer with a banded, repeated weight matrix.
  \item The gradient of a shared weight is the sum over every position where
        it is used.
\end{itemize}

\paragraph{Coming up.}
Section 4.3 adds the remaining building blocks (activation, pooling and the
classifier head) and assembles a complete CNN.

\end{document}
